\section{Relative comparison at an exact physical observation time}
\label{sec:relative-physical-windows}
The previous section gives a denominator for an unchanged, unsmoothed
window. We now compare two different events relative to such a
denominator. The time of physical observation is deterministic. The
initial law in this section is $\nu_R^*$ at a section collision, not
the stationary suspension law. The window widths will be specified;
we do not replace a singleton lattice event by a larger one in the
raw theorem.

\subsection{A three-dimensional denominator from the four-dimensional band}
Put
\begin{equation}\label{eq:physical-window-exponents}
 \rho=\frac{67}{1400},\qquad \sigma=\frac1{25},\qquad
 \epsilon_n=n^{-11/1400}\sqrt{\log(2+n)}.
\end{equation}
For a positive definite $d\times d$ matrix $C$, $g_C$ denotes its
normalized $d$-dimensional Gaussian density.

\begin{lemma}[Projection of the unsmoothed window theorem]
\label{lem:projected-window-denominator}
Let $P_{n,R}:\R^4\to\R^3$ be any linear map which extends to an
invertible $4\times4$ matrix $S_{n,R}$ with
$\|S_{n,R}\|+\|S_{n,R}^{-1}\|\le C_0$. Under $\nu_R^*$, put
$Y=P_{n,R}U_{n,R}/\sqrt n$ and
$\Sigma_{n,R}=P_{n,R}D_RP_{n,R}^{\mathsf T}$.
For $|\xi|\le A$ and half-widths satisfying
$c_0n^{-\sigma}\le h_i\le C_0n^{-\sigma}$, $1\le i\le3$,
\begin{equation}\label{eq:projected-window-denominator}
 \nu_R^*\{Y\in B_3(\xi,\mathbf h)\}
  =\int_{B_3(\xi,\mathbf h)}g_{\Sigma_{n,R}}(y)\dd y
                +O_{A,c_0,C_0}\left(\prod_{i=1}^3(2h_i)\epsilon_n\right).
\end{equation}
The event is the actual projection of the record. No restriction of an
$L^1(\R^4)$ Fourier estimate to a measure-zero frequency plane is used.
\end{lemma}
\begin{proof}
Write $Y'=S_{n,R}U_{n,R}/\sqrt n$ and truncate its fourth coordinate
to $[-T,T]$, where $T=n^{1/400}$. The characteristic error for $Y'$
on the cube $|v_i|\le b$, $b=c(C_0)n^\rho$, has integral at most
$C\mathcal E_n(s_R)$. Indeed $S_{n,R}^{\mathsf T}$ maps this cube
into the retained isotropic ball $|v|\le2n^\rho$, and its determinant
and inverse determinant are bounded. This is a four-dimensional
change of variables, not a trace estimate.

Apply the product envelopes \eqref{eq:box-lower-envelope} to
$B_3(\xi,\mathbf h)\times[-T,T]$. Their $L^1$ norms are
$O(TV_3)$, where $V_3=\prod_{i=1}^3(2h_i)$. Fourier comparison thus
costs $CT V_3\mathcal E_n(s_R)$. We keep a sharper bound for the
Gaussian envelope error. Uniform ellipticity of $S_{n,R}D_RS_{n,R}^{\mathsf T}$
gives a joint density bounded by $Ce^{-c|y|^2}$. For a term containing
an error factor $E_{I_i}$ with $i\le3$, integrate that factor and
the other two output factors in $L^1$; integrate the fourth majorant
against $e^{-cy_4^2}$, using $0\le U_{[-T,T]}\le3$. This is at most
$C V_3/(bh_i)$, independently of $T$. The fourth-coordinate error
term has integral at most $CV_3/b$. The same bounds apply to both
envelopes by \eqref{eq:box-lower-envelope}.

The fourth coordinate of $Y'$ has bounded sixty-fourth moment,
uniformly in $n,R,S_{n,R}$, by
\eqref{eq:all-stopped-moment-bound}. Removing its truncation costs
at most $C T^{-64}$. For the Gaussian, the corresponding cost on
the output box is at most $C V_3 e^{-cT^2}$. Consequently the
absolute probability error divided by $V_3$ is at most
\begin{equation}\label{eq:projection-error-budget}
 C\left[T\mathcal E_n(s_R)+\sum_{i=1}^3(bh_i)^{-1}+b^{-1}
                     +T^{-64}/V_3+e^{-cT^2}\right].
\end{equation}
The four relevant power margins are
\[
 e-\frac1{400}=\frac{23}{2800},\qquad
 \rho-\sigma=\frac{11}{1400},\qquad
 \frac{64}{400}-3\sigma=\frac1{25},\qquad \rho>\frac{11}{1400}.
\]
Also $v_0-1/400>23/2800$, and the norms of $s_R$ are uniform.
Each term in \eqref{eq:projection-error-budget} is therefore
$O(\epsilon_n)$. This proves the lemma, including the discrete endpoint
conventions inherited from the envelopes.
\end{proof}

\subsection{Maximal actual-return increments and inverse clocks}
\begin{lemma}[A deterministic return-window maximal estimate]
\label{lem:actual-return-window-maximal}
For every fixed real $p>2$, every integer $q\ge1$, and $n\ge q$,
\begin{equation}\label{eq:actual-return-window-maximal}
 \int\max_{|j|\le q}|U_{n+j,R}-U_{n,R}|^p\dd\nu_R^*
                       \le C_p q^{p/2}.
\end{equation}
\end{lemma}
\begin{proof}
Every centered return increment over a deterministic interval of
length $l$ has $p$th moment at most $C_p l^{p/2}$, by invariance of
$F_R^*$ and \eqref{eq:all-stopped-moment-bound}. Enlarge $q$ to a
power of two $Q<2q$. At scale $l$, the maximum over its $Q/l$
aligned intervals has $L^p$ norm at most
$C_p Q^{1/p}l^{1/2-1/p}$. A prefix is a sum of at most one interval
per scale. Minkowski's inequality and $p>2$ give a convergent geometric
sum bounded by $C_p\sqrt Q$. Translation by the deterministic index
$n$ is justified by return-map invariance. The backward part uses
suffixes of the length-$q$ interval ending at $n$, and has the same
bound. Adding the two parts proves the assertion. The maximum is over
a short deterministic return window, not over all marks in a
collision Fourier theorem.
\end{proof}

Start the physical orbit at $x\in Y_R^*$, with physical time zero at
that outgoing collision. Let $W_R^Y(t)$ be its lattice displacement
and $C_R^Y(t)$ its completed collision count, using the physical
convention of Section~\ref{sec:uniform-fluid-clock}. Put
\begin{equation}\label{eq:section-start-physical-vector}
 V_R^Y(t)=\left(W_R^Y(t),C_R^Y(t)-t/\bar\tau_R\right),\quad
 B_R(x_1,x_2,x_3,x_4)=(x_1,x_2,x_3-x_4/\bar\tau_R).
\end{equation}
The chosen displacement convention can differ from the sum of completed
flight labels by one uniformly bounded cell-position term; that term
is retained in the estimates below. Define the actual completed-return
clock and its deterministic mean index by
\begin{equation}\label{eq:physical-return-index}
 M_R(t)=\max\{m\ge0:T_{m,R}\le t\},\qquad
 n_R(t)=\left\lfloor c_*t/\bar\tau_R\right\rfloor.
\end{equation}
The positive lower free-flight length makes the maximum finite and
bounds it deterministically by $C(1+t)$.

\begin{proposition}[A coupling at the deterministic physical time]
\label{prop:physical-window-coupling}
For every fixed $P>0$, there is $C_P$ such that, uniformly in $R$ and
all sufficiently large $t$,
\begin{equation}\label{eq:physical-window-coupling}
 \nu_R^*\left\{
   |V_R^Y(t)-B_RU_{n_R(t),R}|>C_P t^{3/8}
          \right\}\le C_Pt^{-P}.
\end{equation}
The same bound holds for the difference
$B_RU_{M_R(t),R}-B_RU_{n_R(t),R}$.
\end{proposition}
\begin{proof}
Put $n=n_R(t)$, $q=\lceil t^{5/8}\rceil$; for large $t$, $q<n$
and $n\asymp t$ uniformly. The return flight-time mean is
$\bar\tau_R/c_*$. The equivalences for the increasing times $T_m$
turn $|M_R(t)-n|>q$ into a deviation of $T_{n+q}$ or $T_{n-q}$
from its mean of size at least $c q$. The rounding error is uniformly
bounded. For every fixed $p$, \eqref{eq:all-stopped-moment-bound}
and Markov's inequality therefore give
\[
 \nu_R^*\{|M_R(t)-n|>q\}\le C_p t^{p/2}q^{-p}
                                      \le C_p t^{-p/8}.
\]
Off this event, \eqref{eq:actual-return-window-maximal} bounds the
probability that $|U_{M_R(t)}-U_n|>t^{3/8}$ by
$C_p q^{p/2}t^{-3p/8}\le C_p t^{-p/16}$. The operator norm of
$B_R$ is uniform, and $B_R\bar G_R=0$.

At the last completed section return, the vector is
$B_RJ_{M_R(t)}=B_RU_{M_R(t)}$. Between this return and time $t$,
the change in displacement, collision count, and elapsed time is
bounded by $C r_R^*\circ(F_R^*)^{M_R(t)}+C$. Uniformly bounded
single flights and labels justify this assertion even when the next
return has not yet occurred. Since $M_R(t)\le C(1+t)$,
invariance and the exponential one-return tail give
\[
 \nu_R^*\left\{\max_{0\le j\le C(1+t)}r_R^*\circ(F_R^*)^j
                                               >t^{3/8}\right\}
       \le C(1+t)e^{-ct^{3/8}}.
\]
This is a union bound, not a conditional independence statement.
Choose a fixed $p>\max\{2,16P\}$ and combine the three events.
The count of moment factors is fixed for each $P$. This proves both
assertions with the actual physical observation time left unchanged.
\end{proof}

\subsection{A relative event theorem with its actual denominator}
Recall the uniformly positive physical covariance
\[
 \mathcal V_R=\bar\tau_R^{-1}B_R\Gamma_RB_R^{\mathsf T}
             =(c_*/\bar\tau_R)B_RD_RB_R^{\mathsf T}.
\]
For $h=t^{-1/25}$ and $|\xi|\le A$, let $B_t=B_3(\xi,(h,h,h))$.
On the same initial probability space $(Y_R^*,\nu_R^*)$ define
\begin{align}
 A_t^{\rm ret}&=\{t^{-1/2}B_RU_{n_R(t),R}\in B_t\},\notag\\
 A_t^{\rm last}&=\{t^{-1/2}B_RU_{M_R(t),R}\in B_t\},\notag\\
 A_t^{\rm phys}&=\{t^{-1/2}V_R^Y(t)\in B_t\}.
 \label{eq:three-physical-window-events}
\end{align}
All three are exact events. The third is measured at deterministic
physical time $t$; it uses integer lattice and collision-count
intervals of physical half-width $t^{23/50}$.

\begin{theorem}[An unsmoothed physical denominator and relative replacement]
\label{thm:relative-physical-window}
Uniformly for $R\in I$, $|\xi|\le A$, and large $t$, all three events
in \eqref{eq:three-physical-window-events} have probability
\begin{equation}\label{eq:physical-window-denominator}
 8t^{-3/25}g_{\mathcal V_R}(\xi)
       \left(1+O_A(t^{-11/1400}\sqrt{\log(2+t)})\right).
\end{equation}
For either $A_t=A_t^{\rm last}$ or $A_t^{\rm phys}$,
\begin{equation}\label{eq:relative-physical-window}
 \frac{\nu_R^*(A_t\mathbin\triangle A_t^{\rm ret})}
                    {\nu_R^*(A_t^{\rm ret})}
       \le C_A t^{-11/1400}\sqrt{\log(2+t)}.
\end{equation}
The total variation distance, defined as the supremum over measurable
sets, between the two conditional initial laws is bounded by the same
right-hand side with a changed constant. Consequently the comparison
holds for every common bounded measurable path statistic, with a
factor twice its supremum norm for the difference of expectations.
\end{theorem}
\begin{proof}
Set $n=n_R(t)$ and take
$P_{n,R}=\sqrt{n/t}\,B_R$ in
Lemma~\ref{lem:projected-window-denominator}. Adjoin the fourth row
$(0,0,0,1)$ to obtain $S_{n,R}$. The ratio $n/t$ is bounded above
and away from zero uniformly, as are $\bar\tau_R$ and its reciprocal;
thus the matrix and its inverse satisfy the lemma's bounds.
Its projected covariance is $(n/t)B_RD_RB_R^{\mathsf T}$, which
differs from $\mathcal V_R$ by $O(t^{-1})$. Uniform ellipticity and
the Gaussian formula make the density error $O_A(t^{-1})$ on central
boxes. Replacing the box Gaussian mass by its center value costs
relative $O_A(h)$. The lemma proves
\eqref{eq:physical-window-denominator} for $A_t^{\rm ret}$,
including a lower bound $c_A t^{-3/25}$.

For the event comparison, use
Proposition~\ref{prop:physical-window-coupling} with $P=2$.
Outside a set of probability $O(t^{-2})$, the standardized vector
changes by at most $d_t=Ct^{-1/8}$. Hence the symmetric difference
is contained in that exceptional set and the event that the return
vector lies in $B_t^+\setminus B_t^-$, where the half-widths are
$h+d_t$ and $h-d_t$. For large $t$ they lie between $h/2$ and $2h$.
Apply the projected-window lemma separately to these two boxes and
subtract their probabilities. The Gaussian boundary volume is at
most $C_A h^3 d_t/h$, with
\[
 d_t/h=O(t^{-17/200}).
\]
The two Fourier-envelope errors have total at most
$C_A h^3 t^{-11/1400}\sqrt{\log(2+t)}$. Thus the symmetric-difference
probability divided by the proved denominator is at most
\[
 C_A\left[t^{-17/200}+t^{-11/1400}\sqrt{\log(2+t)}
                                      +t^{-2+3/25}\right].
\]
The first and third terms have strictly faster power decay, so this proves
\eqref{eq:relative-physical-window}. This argument subtracts masses
of enlarged and contracted boxes; it does not require a separate
Fourier estimate for a thinner unresolved boundary slab.
The difference of event probabilities is bounded by their symmetric
difference, proving \eqref{eq:physical-window-denominator} for the
other events as well.

Finally, if $p=\nu_R^*(A_t^{\rm ret})$ and the symmetric-difference
mass is $\delta\le p/2$, then $\nu_R^*(A_t)\ge p/2$.
Writing the two normalized indicator densities shows directly that
the conditional total variation distance is at most $2\delta/p$.
Integration of any common bounded statistic proves the final assertion.
No event was silently substituted in a numerator or a denominator.
\end{proof}

\paragraph{The still required microscopic comparison.}
Theorem~\ref{thm:relative-physical-window} proves a relative comparison
for different events and a nonzero denominator at deterministic
physical time, but at the displayed growing physical window widths
and under a section-start law. The original raw application also
requires fixed lattice labels, its own exact time constraint, and
its stated initial law. Those are not the events of
\eqref{eq:three-physical-window-events}. In particular neither this
theorem nor the signed average estimate in the preceding section
supplies the full pointwise common remainder, the long-time density
jet budget, or the full weighted fixed-count Fourier complement.
The same raw mixed-density target and all its existing mechanisms
are retained. The new results identify precisely a range in which
actual denominators and relative event errors can already be proved,
without presupposing the microscopic theorem.
